\documentclass[../main.tex]{subfiles}
\begin{document}
\chapter{2018-2019学年微积分（一）（下）期末考试}

\section{单项选择题(每小题 3 分，共 18 分)}
\begin{enumerate}
    \item 设$y_{1}(x),y_{2}(x),y_{3}(x)$都是微分方程$y''+a(x)y'+b(x)y=x^{2}$的解，则以下函数中也是该微分方程的解的是（\qquad）.

          \begin{tasks}[label=\textbf{\Alph*.}, label-width=1.5em, column-sep=1em](2)
              \task $y_{1}+y_{2}+y_{3}$
              \task $\frac{3}{2} y_{1}- y_{2}+\frac{1}{2}y_{3}$
              \task $y_{1}-y_{2}$
              \task $2y_{3}$
          \end{tasks}

    \item 以下函数在原点可微的是（\qquad）.

          \begin{tasks}[label=\textbf{\Alph*.}, label-width=1.5em, column-sep=1em](2)
              \task $z=x^{2}+y^{2}$
              \task $z=\sqrt{x^{2}+y^{2}}$
              \task $z=|x-y|$
              \task $z=\sqrt{|xy|}$
          \end{tasks}

    \item 设$F(x,y)=0$是一条平面光滑曲线，则以下说法中正确的是（\qquad）.

          \begin{tasks}[label=\textbf{\Alph*.}, label-width=1.5em, column-sep=1em](2)
              \task $(F_{x},F_{y})$是该曲线的切向量
              \task $(F_{y},F_{x})$是该曲线的法向量
              \task $(-F_{y},F_{x})$是该曲线的切向量
              \task $(-F_{x},F_{y})$是该曲线的法向量
          \end{tasks}

    \item 设平面区域$D$由$x^{2}+y^{2}\leq1$表示，区域$D_{1}$是$D$在第一象限的部分，则$\iint\limits_{D}(xy+\cos x\sin y)\,\mathrm{d}\sigma=$（\qquad）.

          \begin{tasks}[label=\textbf{\Alph*.}, label-width=1.5em, column-sep=1em](2)
              \task $0$
              \task $4\iint\limits_{D_{1}}xy\,\mathrm{d}\sigma$
              \task $4\iint\limits_{D_{1}}\cos x\sin y\,\mathrm{d}\sigma$
              \task $4\iint\limits_{D_{1}}(xy+\cos x\sin y)\,\mathrm{d}\sigma$
          \end{tasks}

    \item \begin{minipage}[t]{0.93\linewidth}
              设区域$\Omega$由$z=x^{2}+y^{2}$与$z=1$围成，\(I=\iiint\limits_{\Omega}f(z)\,\mathrm{d}v\)，则以下表达式\textbf{错误}的是（\qquad）.\par\smallskip
              \begin{tasks}[label=\textbf{\Alph*.}, label-width=1.5em, column-sep=1em](1)
                  \task $\displaystyle I=\pi\int_{0}^{1}zf(z)\,\mathrm{d}z$ \qquad
                  \task $\displaystyle I=2\pi\int_{0}^{1}r\,\mathrm{d}r\int_{r^{2}}^{1}f(z)\,\mathrm{d}z$
                  \task $\displaystyle I=\int_{-1}^{1}\,\mathrm{d}x\int_{-\sqrt{1-x^{2}}}^{\sqrt{1-x^{2}}}\,\mathrm{d}y\int_{x^{2}+y^{2}}^{1}f(z)\,\mathrm{d}z$
                  \task $\displaystyle I=2\pi\int_{0}^{1}f(z)\,\mathrm{d}z\int_{0}^{z}r\,\mathrm{d}r$
              \end{tasks}
          \end{minipage}

    \item 已知级数$\sum\limits_{n=1}^{\infty}a_{n}^{2}$收敛，则以下说法中正确的是（\qquad）.
          \begin{tasks}[label=\textbf{\Alph*.}, label-width=1.5em, column-sep=1em](2)
              \task $\displaystyle\sum_{n=1}^{\infty}|a_n|$收敛
              \task $\displaystyle\sum_{n=1}^{\infty}a_na_{n+1}$收敛
              \task $\displaystyle\sum_{n=1}^{\infty}a_n$收敛
              \task $\displaystyle\sum_{n=1}^{\infty}(-1)^na_n$收敛
          \end{tasks}
\end{enumerate}

\section{填空题(每小题 4 分，共 16 分)}
\begin{enumerate}
    \setcounter{enumi}{6}
    \item 微分方程$y''+2y'+y=x+2$的通解为\underline{\hspace{5em}}.
          \vspace{1em}

    \item 设$z=f(x^{2}+y^{2},xy)$，其中$f$有连续偏导数，则$\dfrac{\partial z}{\partial x}=$\underline{\hspace{5em}}.
          \vspace{1em}

    \item 设$f(x,y,z)=x+y^{3}+z^{5}$，则$\text{div }\textbf{grad}f=$\underline{\hspace{5em}}.
          \vspace{1em}

    \item $\displaystyle\sum_{n=1}^{\infty}\frac{2^{n}}{n3^{n}}=$\underline{\hspace{5em}}.
          \vspace{1em}
\end{enumerate}

\section{基本计算题(每小题 7 分，共 42 分)}
\begin{enumerate}
    \setcounter{enumi}{10}
    \item 求点$A(1,2,3)$到直线$L:\dfrac{x-6}{-1}=\dfrac{y-1}{1}=\dfrac{z-6}{-2}$的距离.
          \vspace{11em}

    \item 设方程组
          \[
              \begin{cases}
                  x^{2}-2x+y^{2}+\ln z+z^{3}=1, \\
                  x+y+z=1
              \end{cases}
          \]
          在包含点$(0,0,1)$的一个邻域上确定隐函数$y=y(x),z=z(x)$，求\( \left.\frac{\mathrm{d}y}{\mathrm{d}x}\right|_{x=0}\)，\(\left.\frac{\mathrm{d}z}{\mathrm{d}x}\right|_{x=0} \).
          \vspace{8em}

    \item 求函数$f(x,y)=4x+2xy-x^{2}-3y^{2}$的极值.
          \vspace{8em}

    \item 求$I=\iiint\limits_{\Omega}(x+y+z)\,\mathrm{d}x\,\mathrm{d}y\,\mathrm{d}z$，其中$\Omega$是三坐标面与平面$x+y+z=1$所围成的四面体.
          \vspace{9em}

    \item 求曲线积分$I=\int_{L}(2xy-y^{2}\cos x)\,\mathrm{d}x+(x^{2}-2y\sin x)\,\mathrm{d}y$，其中曲线$L$沿抛物线$y=\pi x-x^{2}+1$从$A(0,1)$到$B(\pi,1)$.
          \vspace{9em}

    \item 判断级数$\displaystyle\sum_{n=1}^{\infty}\dfrac{2+(-3)^{n}}{(2n-1)!!}$的敛散性，若收敛指出是绝对收敛还是条件收敛，并说明理由.
          \vspace{9em}
\end{enumerate}

\section{应用题(每小题 7 分，共 14 分)}
\begin{enumerate}
    \setcounter{enumi}{16}
    \item 求曲面$z=x^{2}+y^{2},0\leq z\leq2$的面积.
          \vspace{10em}

    \item 计算曲面积分$I=\oiint\limits_{S}3xz\,\mathrm{d}y\,\mathrm{d}z+yz\,\mathrm{d}z\,\mathrm{d}x-z^{2}\,\mathrm{d}x\,\mathrm{d}y$，其中$S$是曲面$z=\sqrt{x^{2}+y^{2}}$与$z=\sqrt{2-x^{2}-y^{2}}$所围立体的表面外侧.
          \vspace{10em}
\end{enumerate}

\section{综合题(每小题 5 分，共 10 分)}
\begin{enumerate}
    \setcounter{enumi}{18}
    \item 设$L$是圆周曲线$(x-1)^{2}+(y-1)^{2}=1$（取正向），$f(x)$为正值连续函数. 证明：
          \[
              \oint_{L}xf(y)\,\mathrm{d}y-\frac{y}{f(x)}\,\mathrm{d}x\geq2\pi.
          \]
          \vspace{10em}

    \item 证明：当$-\pi<x<\pi$时，$\sum_{n=1}^{\infty}\frac{(-1)^{n-1}\sin nx}{n}=\frac{x}{2}$.
          \vspace{10em}
\end{enumerate}

\chapter{2018-2019学年微积分（一）（下）期末考试参考答案}
\section{单项选择题(每小题 3 分，共 18 分)}
\begin{enumerate}
    \item \textbf{Solution}. B.

          对于非齐次方程$y''+a(x)y'+b(x)y=f(x)$的解$y_1,y_2,y_3$，

          其线性组合$c_1y_1+c_2y_2+c_3y_3$也是该方程的解当且仅当
          \[
              c_1+c_2+c_3=1,
          \]
          四个选项中只有 B 满足
          \[
              \frac32-1+\frac12=1.
          \]

    \item \textbf{Solution}. A.

          对于A选项，$\lim_{(x,y)\to(0,0)}\frac{z(x,y)-z(0,0)-z_x(0,0)x-z_y(0,0)y}{\sqrt{x^2+y^2}}=\lim_{(x,y)\to(0,0)}\sqrt{x^2+y^2}=0$，

          因此$z=x^2+y^2$在原点可微.

          对于B选项，$\lim_{(x,y)\to(0,0)}\frac{z(x,y)-z(0,0)-z_x(0,0)x-z_y(0,0)y}{\sqrt{x^2+y^2}}=\lim_{(x,y)\to(0,0)}\frac{\sqrt{x^2+y^2}}{\sqrt{x^2+y^2}}=1$，

          因此$z=\sqrt{x^2+y^2}$在原点不可微.

          对于C选项，由于$z(x,0)=|x|$，显然$z_x(0,0)$不存在，更不可微.

          对于D选项，$\lim_{(x,y)\to(0,0)}\frac{z(x,y)-z(0,0)-z_x(0,0)x-z_y(0,0)y}{\sqrt{x^2+y^2}}=\lim_{(x,y)\to(0,0)}\frac{\sqrt{|xy|}}{\sqrt{x^2+y^2}}$，

          考察路径$y=x$，$\lim_{x\to0}\frac{\sqrt{|x^2|}}{\sqrt{2x^2}}=\frac1{\sqrt2}$；再考察路径$y=0$，上述极限值为0，

          因此$\lim_{(x,y)\to(0,0)}\frac{\sqrt{|xy|}}{\sqrt{x^2+y^2}}$不存在，$z=\sqrt{|xy|}$在原点不可微.

    \item \textbf{Solution}. C.

          切向量可取作$\left(1,\frac{\mathrm{d}y}{\mathrm{d}x}\right)$，而由隐函数定理可知$\frac{\mathrm{d}y}{\mathrm{d}x}=-\frac{F_x}{F_y}$，

          因此切向量为$\left(1,-\frac{F_x}{F_y}\right)$，亦可取$\left(-F_y,F_x\right)$. C选项正确.

    \item \textbf{Solution}. A.

          已知区域 \(D:x^2+y^2\le1\) 关于 \(x\) 轴、\(y\) 轴均对称.

          由于函数 \(xy\) 关于 \(x\) 或 \(y\) 都是奇函数，而 \(\cos x\sin y\) 关于 \(y\) 是奇函数，所以
          \[
              \iint\limits_D xy\,\mathrm{d}\sigma=0,\qquad
              \iint\limits_D \cos x\sin y\,\mathrm{d}\sigma=0.
          \]
          因此原积分为 \(0\).

    \item \textbf{Solution}. D.

          由柱坐标可知
          \[
              0\le r\le1,\qquad r^2\le z\le1,
          \]
          故
          \[
              I=2\pi\int_0^1 r\,\mathrm{d}r\int_{r^2}^1 f(z)\,\mathrm{d}z.
          \]
          若先对 \(z\) 积分，则 \(0\le z\le1,\ 0\le r\le\sqrt z\)，所以
          \[
              I=2\pi\int_0^1 f(z)\,\mathrm{d}z\int_0^{\sqrt z}r\,\mathrm{d}r
              =\pi\int_0^1 zf(z)\,\mathrm{d}z.
          \]
          因此 A、B 正确，C 是直角坐标写法，也正确；D 中上限应为 \(\sqrt z\)，故错误.

    \item \textbf{Solution}. B.

          对于A、C选项，令$a_n=\frac{1}{n}$，则\(\sum a_n^2\) 收敛，但 \(\sum |a_n|\) 与 \(\sum a_n\) 均发散.

          对于B选项，$|a_na_{n+1}|\le\frac{a_n^2+a_{n+1}^2}{2}$，由比较判别法可知 \(\sum a_na_{n+1}\) 绝对收敛.

          对于D选项，令$a_n=\frac{(-1)^n}{n}$，则 \(\sum a_n^2=\sum \frac{1}{n^2}\) 收敛，但 \(\sum (-1)^na_n = \sum \frac{1}{n}\) 发散.

\end{enumerate}

\section{填空题(每小题 4 分，共 16 分)}
\begin{enumerate}
    \setcounter{enumi}{6}
    \item \textbf{Solution}. $y=(C_1+C_2x)\mathrm{e}^{-x}+x$.

          方程对应的齐次方程$y''+2y'+y=0$的特征方程为
          \[
              r^2+2r+1=(r+1)^2=0,
          \]
          故齐次解为 \((C_1+C_2x)\mathrm{e}^{-x}\). 由于$\lambda=0$不是特征方程的根，所以可设特解为 \(y^*=ax+b\)，代入得
          \[
              ax+b+2a=x+2,
          \]
          所以 \(a=1,\ b=0\).

    \item \textbf{Solution}. $2xf_{1}+yf_{2}$.

          记
          \[
              u=x^2+y^2,\qquad v=xy,
          \]
          则
          \[
              \frac{\partial z}{\partial x}=f_1u_x+f_2v_x=2xf_1+yf_2.
          \]

    \item \textbf{Solution}. $6y+20z^{3}$.

          因为
          \[
              \operatorname{div}\mathbf{grad}f=\Delta f=f_{xx}+f_{yy}+f_{zz},
          \]
          而
          \[
              f_{xx}=0,\qquad f_{yy}=6y,\qquad f_{zz}=20z^3,
          \]
          所以
          \[
              \operatorname{div}\mathbf{grad}f=6y+20z^3.
          \]

    \item \textbf{Solution}. $\ln3$.

          由
          \[
              -\ln(1-t)=\sum_{n=1}^{\infty}\frac{t^n}{n}\qquad(|t|<1),
          \]
          取 \(t=\dfrac23\)，得
          \[
              \sum_{n=1}^{\infty}\frac{2^n}{n3^n}
              =-\ln\left(1-\frac23\right)=\ln3.
          \]
\end{enumerate}

\section{基本计算题(每小题 7 分，共 42 分)}
\begin{enumerate}
    \setcounter{enumi}{10}
    \item \textbf{Solution}. 取直线$L$上一点$B(6,1,6)$，则$\overrightarrow{AB}=\{5,-1,3\}$，直线的方向向量为$\boldsymbol{s}=\{-1,1,-2\}$.

          因此所求距离
          \[
              d=\frac{|\overrightarrow{AB}\times\boldsymbol{s}|}{|\boldsymbol{s}|}
              =\frac{|\,\{-1,7,4\}\,|}{\sqrt6}= \sqrt{11}.
          \]

    \item \textbf{Solution}. 方程组对$x$求导，得
          \[
              \begin{cases}
                  2x-2+2yy'+\dfrac{z'}{z}+3z^{2}z'=0, \\
                  1+y'+z'=0.
              \end{cases}
          \]
          代入点$(0,0,1)$得到
          \[
              \begin{cases}
                  -2+4z'(0)=0, \\
                  1+y'(0)+z'(0)=0.
              \end{cases}
          \]
          解得
          \[
              y'(0)=-\frac32,\qquad z'(0)=\frac12.
          \]

    \item \textbf{Solution}. 令
          \[
              \begin{cases}
                  f_{x}=4+2y-2x=0, \\
                  f_{y}=2x-6y=0.
              \end{cases}
          \]
          解得唯一驻点$(3,1)$.
          又
          \[
              f_{xx}=-2,\quad f_{yy}=-6,\quad f_{xy}=2,\quad \Delta=f_{xx}f_{yy}-f_{xy}^{2}=8>0,
          \]
          故$f(x,y)$在$(3,1)$处取得极大值
          \[
              f(3,1)=6.
          \]

    \item \textbf{Solution}. 由轮换对称性可得
          \[
              I=3\iiint\limits_{\Omega}z\,\mathrm{d}v.
          \]
          用截面法，对于固定的$z$，截面是底边为$(1-z)$的等腰直角三角形，面积为$\frac{(1-z)^{2}}{2}$，因此
          \[
              \begin{aligned}
                  I = {} & 3\int_{0}^{1}z\cdot\frac{(1-z)^{2}}{2}\,\mathrm{d}z \\
                  = {} & \frac{3}{2}\int_{0}^{1}z(1-z)^2\,\mathrm{d}z        \\
                  = {} & \frac{3}{2}\cdot\frac{2!}{4!} = \frac{1}{8}.
              \end{aligned}
          \]

    \item \textbf{Solution}. 设
          \[
              P = 2xy-y^{2}\cos x,\qquad Q = x^{2}-2y\sin x,
          \]
          则
          \[
              P_y=Q_x=2x-2y\cos x,
          \]
          所以积分与路径无关，且
          \[
              \left(2xy-y^{2}\cos x\right)\,\mathrm{d}x+\left(x^{2}-2y\sin x\right)\,\mathrm{d}y=\mathrm{d}(x^{2}y-y^{2}\sin x).
          \]

          所以
          \[
              I = \left(x^{2}y-y^{2}\sin x\right)\Big|_{(0,1)}^{(\pi,1)} = \pi^2.
          \]

    \item \textbf{Solution}. 对正项级数
          \[
              \sum_{n=1}^{\infty}\frac{2}{(2n-1)!!}
          \]
          与
          \[
              \sum_{n=1}^{\infty}\frac{3^{n}}{(2n-1)!!}
          \]
          分别用比值判别法，

          因$\lim_{n\to\infty}\frac{2}{(2n+1)!!}\cdot\frac{(2n-1)!!}{2}=\lim_{n\to\infty}\frac{1}{2n+1}=0<1$，所以$\sum\limits_{n=1}^{\infty}\frac{2}{(2n-1)!!}$收敛；

          因$\lim_{n\to\infty}\frac{3^{n+1}}{(2n+1)!!}\cdot\frac{(2n-1)!!}{3^n}=\lim_{n\to\infty}\frac{3}{2n+1}=0<1$，所以$\sum\limits_{n=1}^{\infty}\frac{3^{n}}{(2n-1)!!}$收敛.

          因此$\sum\limits_{n=1}^{\infty}\frac{2+(-3)^{n}}{(2n-1)!!}$绝对收敛.
\end{enumerate}

\section{应用题(每小题 7 分，共 14 分)}
\begin{enumerate}
    \setcounter{enumi}{16}
    \item \textbf{Solution}. 曲面面积元
          \[
              \mathrm{d}S=\sqrt{1+4x^{2}+4y^{2}}\,\mathrm{d}x\,\mathrm{d}y.
          \]
          投影区域为$D:x^{2}+y^{2}\le2$，故
          \[
              S=\iint\limits_{D}\sqrt{1+4x^{2}+4y^{2}}\,\mathrm{d}x\,\mathrm{d}y
              =2\pi\int_{0}^{\sqrt2}r\sqrt{1+4r^{2}}\,\mathrm{d}r
              =\left.\frac{\pi}{6}\left(1+4r^2\right)^{\frac{3}{2}}\right|_{0}^{\sqrt2}
              =\frac{13\pi}{3}.
          \]

    \item \textbf{Solution}. 由 Gauss 公式，
          \[
              I=\iiint\limits_{\Omega}2z\,\mathrm{d}v.
          \]
          用球坐标代换，可得
          \[
              I=2\int_{0}^{2\pi}\,\mathrm{d}\theta\int_{0}^{\frac{\pi}{4}}\cos\varphi\sin\varphi\,\mathrm{d}\varphi\int_{0}^{\sqrt2}\rho^{3}\,\mathrm{d}\rho=\pi.
          \]
\end{enumerate}

\section{综合题(每小题 5 分，共 10 分)}
\begin{enumerate}
    \setcounter{enumi}{18}
    \item \textbf{Proof}. 由 Green 公式，得
          \[
              \oint\limits_{L}xf(y)\,\mathrm{d}y-\frac{y}{f(x)}\,\mathrm{d}x
              =\iint\limits_{D}\left(f(y)+\frac{1}{f(x)}\right)\,\mathrm{d}x\,\mathrm{d}y,
          \]
          由区域$D:(x-1)^{2}+(y-1)^{2}\le1$的轮换对称性，得
          \[
              \iint\limits_{D}f(y)\,\mathrm{d}x\,\mathrm{d}y=\iint\limits_{D}f(x)\,\mathrm{d}x\,\mathrm{d}y.
          \]
          又因 \(f(x)>0\)，有
          \[
              f(x)+\frac{1}{f(x)}\ge2.
          \]
          从而
          \[
              \begin{aligned}
                  \oint\limits_{L}xf(y)\,\mathrm{d}y-\frac{y}{f(x)}\,\mathrm{d}x
                   & =\iint\limits_D\left(f(x)+\frac1{f(x)}\right)\,\mathrm{d}x\,\mathrm{d}y \\
                   & \ge 2\iint\limits_D\,\mathrm{d}x\,\mathrm{d}y=2\pi.
              \end{aligned}
          \]

    \item \textbf{Proof}. 将函数$f(x)=\dfrac{x}{2}$在$(-\pi,\pi)$上作$2\pi$周期延拓并展开为傅立叶级数.
          由奇偶性知$a_{n}=0$，而
          \[
              b_{n}=\frac{2}{\pi}\int_{0}^{\pi}\frac{x}{2}\sin nx\,\mathrm{d}x=-\frac{1}{n\pi}\int_{0}^{\pi}x\,\mathrm{d}\cos nx = \frac{(-1)^{n-1}}{n}.
          \]
          因此其正弦级数为
          \[
              \sum_{n=1}^{\infty}\frac{(-1)^{n-1}\sin nx}{n},
          \]
          由Dirichlet收敛定理可得当$-\pi<x<\pi$时，
          \[
              \sum_{n=1}^{\infty}\frac{(-1)^{n-1}\sin nx}{n}=\frac{x}{2}.
          \]
\end{enumerate}
\end{document}





